\section{What the rules guarantee}\label{sec:properties}

The rules of Definition~\ref{def:rules} are local: each one looks at a single entry and its immediate premises. What one wants to know
about a long record is global: whether a conclusion fifty steps down a chain still carries every condition and every input from the top.
This section shows that the local rules are enough. The proofs are short, and that is the point: the rules are precise enough to be
checked by a compiler and by a proof assistant alike. Appendix~\ref{app:lean} lists the corresponding machine-checked statements.

\subsection{Checking a record in one pass}

We model a checker that reads a record line by line. Its state $S$ is a finite map from keys to the stored information
$S(k)=(\mathit{frame},\mathit{status},\mathit{inputs},\mathit{prem})$, together with the declared maps. Reading a line either fails or
extends the state:
\begin{itemize}[leftmargin=1.5em]
\item an \emph{input} $k$ is accepted if $k$ is new, and stored with $\mathit{inputs}=\{k\}$ and no premises;
\item a \emph{claim} or \emph{thesis} $k$ is accepted if $k$ is new, and stored with no inputs and no premises;
\item a \emph{derivation} $k$ with premises $p_1,\dots,p_m$ ($m\ge1$) is accepted if $k$ is new, every $p_j$ is stored, lies in the
frame of $k$ or in $*$, has status at least that of $k$, and the declared inputs of $k$ are exactly the union of the inputs of the $p_j$;
\item a \emph{conversion} $k=\mathrm{convert}(g,p)$ is accepted if $k$ is new, $g\colon A\to B$ is declared, $p$ is stored in frame $A$,
$k$ is in frame $B$, $k$ is not stronger than $p$, and the inputs of $k$ are those of $p$;
\item a \emph{comparison} or \emph{declaration of independence} is accepted if (R4), respectively (R3), holds for it; it does not change the state;
\item a \emph{map declaration} is accepted if the name of the map is new.
\end{itemize}
A record is \emph{accepted} if every line is accepted in turn, starting from the empty state. This is exactly the procedure of the
Lean checker of Section~\ref{sec:latex}, and it is the reason for rule (R1): because nothing may be used before it is declared, a single
pass suffices.

\subsection{An invariant of accepted records}

\begin{definition}[Good state]\label{def:good}
A state $S$ is \emph{good} if for every stored key $k$ with $S(k)=(f,s,I,P)$:
\begin{enumerate}[label=(\alph*),nosep,leftmargin=2em]
\item every premise $p\in P$ is stored, with $S(p)=(f',s',I',P')$ such that $s\le s'$ and $I'\subseteq I$;
\item every $x\in I$ either is $k$ itself and $P$ is empty (that is, $k$ is an input), or lies in the inputs of some premise $p\in P$.
\end{enumerate}
\end{definition}

\begin{theorem}[Accepted records are good]\label{thm:good}
The state reached after an accepted record is good.
\end{theorem}

\begin{proof}
By induction on the number of lines. The empty state is good because it stores nothing. Suppose $S$ is good and a line is accepted,
giving $S'$. Comparisons, declarations of independence and map declarations leave the stored keys unchanged, so $S'$ is good. Otherwise
exactly one new key $k$ is added, and since $k$ was new, no earlier entry changes; the earlier entries still satisfy (a) and (b) in $S'$
because their premises are still stored with the same information. For the new key, (a) and (b) are exactly what the acceptance test
checked: an input has $I=\{k\}$ and no premises; a claim has neither; a derivation or conversion has stored premises of at least its
status whose inputs are contained in $I$, and every element of $I$ comes from some premise.
\end{proof}

\subsection{No promotion, at any distance}

Say that $a$ is an \emph{ancestor} of $k$ if $a$ can be reached from $k$ by following premises one or more times.

\begin{theorem}[No promotion]\label{thm:nopromotion}
In an accepted record, if $a$ is an ancestor of $k$ then $\mathit{status}(k)\le\mathit{status}(a)$.
\end{theorem}

\begin{proof}
Induction on the length of the path from $k$ to $a$. For a direct premise this is part (a) of Theorem~\ref{thm:good}. For a longer path
$k\to m\to\dots\to a$, the induction hypothesis gives $\mathit{status}(m)\le\mathit{status}(a)$, part (a) gives
$\mathit{status}(k)\le\mathit{status}(m)$, and the order on statuses is transitive.
\end{proof}

\begin{corollary}[Proved rests only on proved]\label{cor:proved}
In an accepted record, every ancestor of an entry with status \st{proved} has status \st{proved}. Equivalently, an entry with an ancestor
below \st{proved} is not \st{proved}.
\end{corollary}

\begin{proof}
\st{proved} is the top of the order, so $\st{proved}\le\mathit{status}(a)$ forces $\mathit{status}(a)=\st{proved}$.
\end{proof}

For the example: the narrowed prediction~\eqref{eq:interval} has the conditional preparation as an ancestor, two steps up. By
Theorem~\ref{thm:nopromotion} it can be at most \st{conditional}, however many further conclusions are later drawn from it.

\subsection{Inputs are never lost and never invented}

\begin{theorem}[Inputs]\label{thm:inputs}
In an accepted record:
\begin{enumerate}[label=(\roman*),nosep,leftmargin=2em]
\item if $a$ is an ancestor of $k$, then $\mathit{inputs}(a)\subseteq\mathit{inputs}(k)$;
\item every input of $k$ is $k$ itself (if $k$ is an input) or an input of one of its premises;
\item if $\mathit{inputs}(k)=\emptyset$, then no ancestor of $k$ has an input, and in particular no ancestor is an input entry.
\end{enumerate}
\end{theorem}

\begin{proof}
(i) Induction on the path, using part (a) of Theorem~\ref{thm:good} and transitivity of $\subseteq$. (ii) is part (b) of
Theorem~\ref{thm:good}. (iii) follows from (i): an ancestor with a non-empty input set would force $\mathit{inputs}(k)\neq\emptyset$;
and an input entry $a$ has $\mathit{inputs}(a)=\{a\}\neq\emptyset$.
\end{proof}

Part (iii) is what makes the declaration of independence meaningful. If a record declares that a result is independent of all inputs and
the record is accepted, then no measured or fitted value occurs anywhere above that result. In the example, the measurement is an
ancestor of the narrowed prediction, so the prediction carries it in its input set, and a declaration that the prediction is
independent of inputs is rejected.

\subsection{Frames change only through declared maps}

To state the frame property we look at a single conclusion and the tree of steps that produced it. A \emph{derivation tree} has
leaves (entries without premises) and three kinds of inner node: a step from one premise, a join of two premises, and a conversion
through a map $g$ into a frame $b$. A tree is \emph{well formed} relative to a set $D$ of declared maps if no node is stronger than
its children, the two children of a join lie in the same frame, and every conversion uses a map in $D$ whose source is the frame of its
child.

\begin{theorem}[Frames]\label{thm:frames}
In a well-formed tree, the frame of every leaf reaches the frame of the root through a chain of declared maps. If no maps are
declared, every leaf lies in the frame of the root.
\end{theorem}
\noindent The tree model has no convention-free frame; leaves in $*$, which the record checker admits as premises of any frame, are
excluded here. The frame property is proved for trees; the record checker enforces the corresponding local condition on every line.


\begin{proof}
Induction on the tree. A leaf reaches its own frame through the empty chain. A step keeps the frame of its child. For a join, both
children have the same frame, which is the frame of the join, and the induction hypothesis applies to each side. For a conversion
through $g\colon a\to b$, every leaf of the child reaches $a$ by induction, and appending $g$ gives a chain to $b$. If $D$ is empty, the
only chains are empty, so every leaf lies in the frame of the root.
\end{proof}

Theorem~\ref{thm:frames} is the formal content of rule (R4): there is no way to arrive in the Boosted frame from a Lab statement except
through $\mathit{lorentz}_z$ or another declared map.

\subsection{The hash chain binds conclusions to premises}

The checker also reads the hash and the list of source hashes recorded for every entry. It accepts the hash part of a line if the
entry's hash has not occurred before and, for a derivation or conversion, if the recorded sources are exactly the hashes recorded
earlier for its premises, in order; leaves must have no sources.

\begin{proposition}[Hash links]\label{prop:links}
In a record accepted by the hash check, the sources of every derivation and conversion are the hashes of its premises as they were
recorded, and no hash occurs twice. If, in addition, the hashes are computed as in Definition~\ref{def:hash} and no MD5 collision
occurs among the entries involved, then changing any field of an entry changes its hash and the hash of every entry of which it is an
ancestor, and leaves the hash of every other entry unchanged.
\end{proposition}

\begin{proof}
The first sentence is what the check tests, line by line. For the second, an entry's hash is computed from its own fields and the
hashes of its premises. Changing a field changes the hashed string and therefore, absent a collision, the hash. By induction along the
paths from the changed entry, every descendant receives a different source list and hence a different hash. An entry that is not a
descendant has the same fields and, by induction on the record, the same premise hashes, and hence the same hash.
\end{proof}

The proposition separates what is proved from what is assumed. That the links are consistent is checked. That a changed statement
yields a changed hash is a property of the hash function, which is assumed and not proved; MD5 is chosen to detect accidental change
and gives no protection against deliberate collisions.

\subsection{What the proofs do not say}

All results above are relative to what is declared. A record that declares a premise-free claim as \st{proved} is accepted, and so are
its consequences. A checked record shows that the stated statuses, inputs and frames are mutually consistent. It does not show that
any statement is true.
